
\documentclass[12pt, twoside]{book}
\usepackage[utf8]{inputenc}
\usepackage[spanish, es-noquoting]{babel}
\usepackage[width=16cm,height=21cm]{geometry}
\usepackage{tablefootnote}
\usepackage{fancyhdr}
\usepackage[backend=biber, style=authoryear, natbib=true]{biblatex}
\usepackage{tabularray}
\usepackage[nottoc]{tocbibind}
\usepackage{hyperref}
\usepackage{amsmath, amsfonts, amssymb, amsthm}
\usepackage{nicematrix}
\pagestyle{fancy}
\setlength{\headheight}{20pt}
\linespread{1.1}
\fancyhead[E]{\thepage\hfill \textit{\nouppercase{\rightmark}}}
\fancyhead[O]{\textit{\nouppercase{\leftmark}}\hfill\thepage}
\renewcommand{\headrulewidth}{0.5pt}
\cfoot{}
\usepackage{amsmath,amssymb,mathtools}
\usepackage{url}
\usepackage{tikz}
\usetikzlibrary{shapes.geometric}
\usepackage[intoc]{nomencl}
\makenomenclature
\renewcommand{\nomname}{Lista de símbolos}
\usepackage{cases}
\usepackage{xparse}
\usepackage{mathtools}
\usepackage{algorithm}
\usepackage{algpseudocode}

\floatname{algorithm}{Algoritmo}
\renewcommand{\listalgorithmname}{Índice de algoritmos}

% 2. Traducir los comandos internos
\algrenewcommand\algorithmicrequire{\textbf{Entrada:}}
\algrenewcommand\algorithmicensure{\textbf{Salida:}}
\algrenewcommand\algorithmicwhile{\textbf{mientras}}
\algrenewcommand\algorithmicdo{\textbf{hacer}}
\algrenewcommand\algorithmicend{\textbf{fin}}
\algrenewcommand\algorithmicif{\textbf{si}}
\algrenewcommand\algorithmicelse{\textbf{si no}}
\algrenewcommand\algorithmicthen{\textbf{entonces}}
\algrenewcommand\algorithmicfor{\textbf{para}}
\algrenewcommand\algorithmicforall{\textbf{para todo}}
\algrenewcommand\algorithmicrepeat{\textbf{repetir}}
\algrenewcommand\algorithmicuntil{\textbf{hasta que}}
\algrenewcommand\algorithmicreturn{\textbf{regresar}}

% Paquete para gestión de cambios
% La opción 'draft' muestra los cambios. 
% La opción 'final' acepta todos los cambios y oculta las marcas (limpia el PDF).
\usepackage[draft, commentsnumbered, trackchanges]{changes}

% Personalización de alias para que coincidan con los botones de TeXStudio
\newcommand{\add}[1]{\added{#1}}
\newcommand{\remove}[1]{\deleted{#1}}
\newcommand{\replace}[2]{\replaced{#1}{#2}} % Ojo: \replaced{nuevo}{viejo}
\newcommand{\alert}[1]{\textbf{\textcolor{red}{#1}}} % Definición manual simple

%=====================================
%
%\newtheorem{definition}{Definición}[chapter]
%\newtheorem{theorem}{Teorema}[chapter]
%\newtheorem{lemma}[theorem]{Lema}
%\newtheorem{coll}[theorem]{Corolario}
%\newtheorem{notation}{Notación}[chapter]
%
%

\usepackage[framemethod=tikz]{mdframed} % Alternativa estable a tcolorbox

% 1. Definimos tu estilo de texto (cuerpo normal, título en negrita)
\newtheoremstyle{estiloTesisPlano}
{0pt}         % Espacio arriba (lo controla mdframed)
{0pt}         % Espacio abajo (lo controla mdframed)
{\normalfont} % Cuerpo del texto sin cursivas
{}            % Sangría
{\bfseries}   % Fuente del encabezado
{.}           % Puntuación después del encabezado
{ }           % Espacio después del encabezado
{}            % Especificación del encabezado

\theoremstyle{estiloTesisPlano}

% 2. Configuramos el diseño visual global de la línea lateral
\mdfsetup{
	hidealllines=true,          % Quita el marco completo
	leftline=true,              % Activa solo la línea izquierda
	linecolor=gray!60,          % Color de la línea
	linewidth=2pt,              % Grosor de la línea
	innerleftmargin=10pt,       % Espacio entre línea y texto
	innerrightmargin=0pt,
	innertopmargin=0pt,
	innerbottommargin=0pt,
	skipabove=2ex,              % Espacio antes del teorema
	skipbelow=2ex               % Espacio después del teorema
}

% 3. Declaramos los entornos (esto une amsthm y mdframed en un solo paso)
\newmdtheoremenv{definition}{Definición}[chapter]
\newmdtheoremenv{theorem}{Teorema}[chapter]
\newmdtheoremenv{lemma}[theorem]{Lema}
\newmdtheoremenv{coll}[theorem]{Corolario}
\newmdtheoremenv{notation}{Notación}[chapter]
